\documentclass{article}

\usepackage{amsmath,amssymb}
\usepackage{xurl}
\usepackage[hidelinks]{hyperref}

% Shared commands for notes in docs/paper-gaps/.

\DeclareUnicodeCharacter{2080}{\ifmmode{}_{0}\else\textsubscript{0}\fi}
\DeclareUnicodeCharacter{2081}{\ifmmode{}_{1}\else\textsubscript{1}\fi}
\DeclareUnicodeCharacter{2082}{\ifmmode{}_{2}\else\textsubscript{2}\fi}
\DeclareUnicodeCharacter{2099}{\ifmmode{}_{n}\else\textsubscript{n}\fi}

\newcommand{\Lean}{\textsc{Lean}}
\ExplSyntaxOn
\NewDocumentCommand{\leanid}{m}{
  \ifmmode
    \textsf{\detokenize{#1}}
  \else
    \group_begin:
    \urlstyle{sf}
    \tl_set:Nn \l_tmpa_tl {#1}
    \tl_replace_all:Nnn \l_tmpa_tl {\_} {_}
    \exp_args:NV \path \l_tmpa_tl
    \group_end:
  \fi
}
\ExplSyntaxOff
\providecommand{\tr}{\operatorname{tr}}
\newcommand{\tnleanIssueBase}{https://github.com/LionSR/TNLean/issues/}
\newcommand{\tnleanPrBase}{https://github.com/LionSR/TNLean/pull/}
\newcommand{\tnleanDocBase}{https://sirui-lu.com/TNLean/docs/}
\newcommand{\ghissue}[1]{\href{\tnleanIssueBase#1}{GitHub issue \##1}}
\newcommand{\ghpr}[1]{\href{\tnleanPrBase#1}{PR \##1}}
\newcommand{\doclink}[1]{\href{\tnleanDocBase#1.html}{\path{#1}}}

% Dirac notation and the identity operator, as in blueprint/src/macros/common.tex.
% \providecommand keeps note-local definitions authoritative.
\usepackage{dsfont}
\providecommand{\ket}[1]{|#1\rangle}
\providecommand{\bra}[1]{\langle #1|}
\providecommand{\braket}[2]{\langle #1 | #2 \rangle}
\providecommand{\ketbra}[2]{|#1\rangle\!\langle #2|}
\providecommand{\Id}{\mathds{1}}
\providecommand{\one}{\mathds{1}}
\providecommand{\C}{\mathbb{C}}
\providecommand{\MN}[1]{M_{#1}(\C)}
\providecommand{\MD}{M_D(\C)}

% External referencing for paper sources.  The build helper writes sanitized
% auxiliary files under build/paper-aux/ and excludes duplicated source labels.
\usepackage{xr}

% Import a generated paper auxiliary file with a caller-supplied label prefix.
% The candidate paths support builds from docs/paper-gaps/, the repository root,
% and build/paper-aux/ itself.
\newcommand{\importpaperaux}[2]{%
  \IfFileExists{../../build/paper-aux/#2.aux}{%
    \externaldocument[#1]{../../build/paper-aux/#2}%
  }{%
    \IfFileExists{build/paper-aux/#2.aux}{%
      \externaldocument[#1]{build/paper-aux/#2}%
    }{%
      \IfFileExists{#2.aux}{%
        \externaldocument[#1]{#2}%
      }{}%
    }%
  }%
}

% General external reference macros
\newcommand{\paperref}[2]{\ref{#1-#2}}
\newcommand{\papereqref}[2]{(\ref{#1-#2})}

% CPSV16 source (1606.00608 / MPDO-22-12-17-2.tex) external document and shortcuts
\importpaperaux{cpsv16-}{1606.00608/MPDO-22-12-17-2-tnlean}
\newcommand{\cpsvref}[1]{\paperref{cpsv16}{#1}}
\newcommand{\cpsveqref}[1]{\papereqref{cpsv16}{#1}}

% Verdict marker: \gapnote{<kind>}{<status>}. Kind is the policy
% classification (clarification, local-correction, scope-restriction,
% unfaithful, false-source, open-gap); status is open, wip (elimination
% underway), resolved, or historical. Machine-read by the site index;
% prints a verdict line.
\newcommand{\gapnote}[2]{\par\noindent\textbf{Verdict:} #1 (#2).\par}


\title{Ising String-Net Operators: the Normalization of the $\sigma$ Tensor}
\author{}
\date{2026-09-24, updated 2026-09-27}

\begin{document}
\maketitle
\gapnote{local-correction}{resolved}

\section{The source statement}

Appendix~D.2, ``Ising string-net'', of Ref.~\cite{Bultinck2017Anyons}
lists the labels $1,\sigma,\psi$ of the Ising category, the fusion rules with
the single nontrivial rule $\sigma\times\sigma=1+\psi$, the quantum dimensions
$d_1=d_\psi=1$, $d_\sigma=\sqrt2$, and the nontrivial $F$-symbols
$F^{\sigma\sigma\sigma}_{\sigma ab}=\pm1/\sqrt2$ and
$F^{\psi\sigma\psi}_{\sigma\sigma\sigma}=F^{\sigma\psi\sigma}_{\psi\sigma\sigma}=-1$.
It then says that the matrix product operator tensors are built from the
$G$-symbols $G^{abc}_{def}=F^{abc}_{def}/(v_ev_f)$, $v_i=\sqrt{d_i}$,
``similarly as for the Fibonacci model''. For the Fibonacci model the source
states that the operator blocks satisfy the fusion rules; for the Ising model
it prints only the fusion rules of the category.\footnote{Source traceability:
    \path{References/1511.08090/AnyonsPEPS.tex}, lines 1257--1268 and
    1305--1323.}
The Ising analogue of the Fibonacci statement is, for the three blocks
$O_1$, $O_\psi$, $O_\sigma$ on a ring of $N$ sites, the algebra
\begin{align}
    O_\psi^2 = O_1,\qquad O_\sigma O_\psi = O_\psi O_\sigma = O_\sigma,\qquad
    O_\sigma^2 = O_1 + O_\psi,
    \label{eq:bmwshv17_ising_fusion_algebra}
\end{align}
with $O_1$ the unit on the admissible closed paths. The same abstract algebra
is printed in Ref.~\cite{Aasen2016IsingDefects} for a different realization,
the defect operators of the Ising lattice model acting on spin and dual-spin
configurations.%
\footnote{Source traceability:
    \path{References/1601.07185/source/Ising-Defects.tex}, lines 1051--1055.}

\section{The chosen normalization}

The formal tensors have physical index the fusion-tree label $(x',\rho,x)$
with $x\in x'\otimes\rho$ (ten labels) and bond index the pair of outer
labels, and their letters are the $F$-symbols
$[F^{a x'\rho}_y]_{y'x}$ themselves, without the factors $v_ev_f$ of the
$G$-symbols. In addition the $\sigma$ tensor is stored as $\sqrt2A_\sigma$,
so that its entries are $0,\pm1,\pm\sqrt2$ rather than
$0,\pm1/\sqrt2,\pm1$, and every entry of the three tensors lies in
$\mathbb Z[\sqrt2]$. The factor $\sqrt2$ is a per-site scalar, so the
periodic operator of the stored tensor is $2^{N/2}$ times that of
$(\sqrt2)^{-1}(\sqrt2A_\sigma)$, and the rule $O_\sigma^2=O_1+O_\psi$ for
the latter is equivalent to
\begin{align}
    O_N(\sqrt2A_\sigma)^2 = 2^N\bigl(O_N(A_1)+O_N(A_\psi)\bigr).
    \label{eq:bmwshv17_ising_scaled_rule}
\end{align}
Both forms are proved.

The omission of the factors $v_ev_f$ does not change the periodic
operators; this is proved in the next section. The identification of the
formal letters with the $F$-symbols was first recorded in the verification
data \path{Notes/OpenProblemsTN/checks/asym_ising_action_data.md}; it is now
a formal statement.

\section{Relation to the source's tensor}

The source's operator tensor (the Fibonacci figure, reused for Ising) is the
crossing of the operator line $f$ with a vertical edge line $e$, with corner
plaquettes $b$ (upper left), $a$ (upper right), $c$ (lower left) and $d$
(lower right), and entry $G^{abc}_{def}\sqrt{v_av_bv_cv_d}$.%
\footnote{Source traceability:
    \path{References/1511.08090/AnyonsPEPS.tex}, lines 1044--1050 and
    1262--1266, figure file \path{StringNetMPO_RHS.pdf}.}
On the fusion-tree labels of the formal tensors the upper label is
$(b,e,a)$, the lower label is $(c,e,d)$, and the bond letters are the
admissible pairs $(b,c)$ and $(a,d)$ of plaquettes above and below the
operator line. With the $F$-symbols and selection rule
$\delta_{abe}\delta_{cde}\delta_{adf}\delta_{bcf}$ of the source, the letters
of $A_1$, $A_\psi$ and $A_\sigma=(\sqrt2)^{-1}\widehat A_\sigma$ are the
$F$-symbols $F^{fce}_{abd}$, and the Ising $G$-symbols have the tetrahedral
symmetry $G^{abc}_{def}=G^{fce}_{abd}$. Hence the nonzero entries of the
source's tensor are those of the formal block multiplied by
$\sqrt{v_av_bv_cv_d}/(v_bv_d)=g(a,d)/g(b,c)$, with $g(u,l)=(v_u/v_l)^{1/2}$,
and the diagonal bond matrix $g$ conjugates every letter of the source's
tensor into the letter of $A_1$, $A_\psi$ or $A_\sigma$:
\begin{align}
    g\,S_f^{kh}\,g^{-1} = A_f^{kh}.
    \label{eq:bmwshv17_ising_gsymbol_conj}
\end{align}
The similarity cancels in the trace, so the periodic operators of the
source's tensors equal $O_N(A_1)$, $O_N(A_\psi)$ and
$O_N(A_\sigma)=2^{-N/2}O_N(\widehat A_\sigma)$ at every length, and the
algebra~(\ref{eq:bmwshv17_ising_fusion_algebra}) holds for the source's
tensors themselves.\footnote{Results:
    \leanid{IsingTwist.isingGSymbol_tetrahedral},
    \leanid{IsingTwist.isingStringNetTensor_conj},
    \leanid{IsingTwist.mpo_isingStringNetTensor},
    \leanid{IsingTwist.mpo_isingStringNetTensor_sigma},
    \leanid{IsingTwist.isMPOFusionAlgebra_isingStringNetTensor}.}

\section{Formal boundary}

The closed-path kernels of $O_1$ and $O_\psi$, the eight products with at
most one factor $\sigma$, and both forms of
(\ref{eq:bmwshv17_ising_scaled_rule}) are proved at every positive length
for the $v$-free tensors, and transfer to the source's $G$-symbol tensors
by~(\ref{eq:bmwshv17_ising_gsymbol_conj}). What remains is the choice of
storing $\widehat A_\sigma=\sqrt2A_\sigma$ so that every entry lies in
$\mathbb Z[\sqrt2]$, a per-site scalar recorded above; nothing is open.

\bibliographystyle{alpha}
\bibliography{references}
\end{document}
